\chapter{Kontrol Optimal Kuadratik dan Estimasi Optimal}
\label{bab:lqr}
\enchap{Quadratic Optimal Control and Optimal Estimation}

\begin{tujuan}
	Setelah mempelajari bab ini, Anda mampu merumuskan indeks kinerja kuadratik, memilih bobot $Q$ dan $R$, menghitung penguatan LQR dari persamaan Riccati diskrit, menerapkan filter Kalman diskrit, serta menjelaskan LQG dan prinsip dasar MPC.
\end{tujuan}

\section{\textit{Linear Quadratic Regulator}}
Penempatan kutub menuntut perancang menentukan lokasi kutub, yang sulit untuk sistem orde tinggi atau banyak masukan. LQR mengganti pertanyaan ``di mana kutubnya'' dengan ``berapa biayanya''. Cari $u[k]$ yang meminimalkan
\begin{equation}
	J = \sum_{k=0}^{\infty}\Big(\mathbf{x}[k]^\top Q\,\mathbf{x}[k] + u[k]^\top R\,u[k]\Big),
\end{equation}
dengan $Q \succeq 0$ menghukum penyimpangan keadaan dan $R \succ 0$ menghukum usaha kontrol.

\paragraph{Solusi horison tak hingga.} Solusi optimalnya adalah umpan balik keadaan konstan $u = -K\mathbf{x}$ dengan
\begin{equation}
	K = (R + \Gamma^\top P\Gamma)^{-1}\Gamma^\top P\Phi,
\end{equation}
dan $P$ solusi \term{persamaan Riccati aljabar diskrit}{discrete algebraic Riccati equation, DARE}:
\begin{equation}
	P = Q + \Phi^\top P\Phi - \Phi^\top P\Gamma(R + \Gamma^\top P\Gamma)^{-1}\Gamma^\top P\Phi.
\end{equation}
Di MATLAB tersedia \texttt{dlqr(Phi,Gam,Q,R)} dan di Python \texttt{control.dlqr}.

\paragraph{Horison berhingga.} Untuk $J_N = \mathbf{x}[N]^\top S\mathbf{x}[N] + \sum_{k=0}^{N-1}(\mathbf{x}^\top Q\mathbf{x} + u^\top Ru)$, pemrograman dinamis memberi rekursi mundur dari $P_N = S$:
\begin{equation}
	K_k = (R + \Gamma^\top P_{k+1}\Gamma)^{-1}\Gamma^\top P_{k+1}\Phi,\qquad P_k = Q + \Phi^\top P_{k+1}(\Phi - \Gamma K_k).
\end{equation}
Bila $N \to \infty$, $P_k$ konvergen ke solusi DARE dan $K_k$ ke penguatan konstan.

\paragraph{Sifat.} Kalang tertutup LQR dijamin stabil bila $(\Phi, \Gamma)$ terkendali dan $(\Phi, Q^{1/2})$ teramati. Perancang hanya memilih bobot dan kutub mengikuti secara otomatis, serta metode ini langsung berlaku untuk banyak masukan. Margin kestabilan LQR kontinu yang terkenal (PM $\ge 60^\circ$) tidak dijamin pada versi diskrit, sehingga margin tetap perlu diperiksa.

\section{Memilih bobot}
Aturan Bryson menormalisasi bobot dengan nilai maksimum yang dapat diterima:
\begin{equation}
	Q_{ii} = \frac{1}{x_{i,\max}^2},\qquad R_{jj} = \frac{1}{u_{j,\max}^2},
\end{equation}
lalu ditala dengan satu faktor $\rho$ pada $R$. Hanya rasio $Q/R$ yang menentukan hasil: $R$ kecil (kontrol ``murah'') memberi respons cepat dengan $u$ besar, sedangkan $R$ besar memberi respons lambat dengan $u$ kecil. Untuk motor posisi dengan $\theta_{\max} = 0{,}1$ rad dan $u_{\max} = 1$, diperoleh $Q = \text{diag}(100, 0)$ dan $R = 1$, setara dengan $Q = \text{diag}(1, 0)$ dan $R = 0{,}01$.

\begin{contoh}[LQR motor posisi]
	Dengan $Q = \text{diag}(1, 0)$, $T = 0{,}1$\,s, dan referensi melalui $\bar N$:
	\begin{center}
		\renewcommand{\arraystretch}{1.2}
		\begin{tabular}{@{}l c c c@{}}
			\toprule
			$R$ & 0,01 & 0,1 & 1 \\
			\midrule
			$K$ & $[7{,}94\ \ 2{,}05]$ & $[2{,}85\ \ 0{,}98]$ & $[0{,}96\ \ 0{,}41]$ \\
			$|z_{cl}|$ & 0,72 & 0,82 & 0,87 \\
			$M_p$ & 3,1\% & 1,3\% & 0\% \\
			$t_s$ (s) & 1,4 & 1,6 & 4,0 \\
			$u_{\max}$ & 7,9 & 2,9 & 1,0 \\
			\bottomrule
		\end{tabular}
	\end{center}
	Satu parameter $R$ menggeser kompromi kecepatan--usaha secara mulus, dengan redaman yang selalu baik (Gambar~\ref{fig:b13-lqr}).
\end{contoh}

\begin{figure}[H]
	\centering
	\begin{tikzpicture}
		\begin{groupplot}[group style={group size=2 by 1,horizontal sep=1.5cm},skd,width=7cm,height=4.8cm,xmin=0,xmax=4,xlabel={$t$ (s)}]
			\nextgroupplot[ylabel={$\theta$},ymin=0,ymax=1.15,legend style={at={(0.98,0.05)},anchor=south east}]
			\addplot[skdgray,dashed,thin,domain=0:4,forget plot] {1};
			\addplot[skdred,const plot] table[x=t,y=y1] {\dat{p14_lqr.dat}};
			\addplot[skdorange,const plot] table[x=t,y=y2] {\dat{p14_lqr.dat}};
			\addplot[skdblue,const plot] table[x=t,y=y3] {\dat{p14_lqr.dat}};
			\legend{$R=0{,}01$, $R=0{,}1$, $R=1$}
			\nextgroupplot[ylabel={$u$},ymin=-1,ymax=8.5]
			\addplot[skdred,const plot] table[x=t,y=u1] {\dat{p14_lqr.dat}};
			\addplot[skdorange,const plot] table[x=t,y=u2] {\dat{p14_lqr.dat}};
			\addplot[skdblue,const plot] table[x=t,y=u3] {\dat{p14_lqr.dat}};
		\end{groupplot}
	\end{tikzpicture}
	\caption{Pengaruh bobot $R$ pada LQR motor posisi}
	\label{fig:b13-lqr}
\end{figure}

\section{Filter Kalman}
Model dengan derau proses dan derau pengukuran adalah
\begin{align}
	\mathbf{x}[k+1] &= \Phi\mathbf{x}[k] + \Gamma u[k] + \Gamma_ww[k],\\
	y[k] &= C\mathbf{x}[k] + v[k],
\end{align}
dengan $w$ dan $v$ derau putih berrerata nol, saling bebas, dan berkovarians $Q_n$ dan $R_n$. Filter Kalman memberi estimasi yang meminimalkan $E\big[\|\mathbf{x} - \hat{\mathbf{x}}\|^2\big]$. Strukturnya sama dengan pengamat \textit{current} (Bab~\ref{bab:obs}), tetapi penguatannya dipilih secara optimal dari $Q_n$ dan $R_n$:

\paragraph{Koreksi.}
\begin{align}
	L_k &= P_k^-C^\top\big(CP_k^-C^\top + R_n\big)^{-1},\\
	\hat{\mathbf{x}}_k &= \hat{\mathbf{x}}_k^- + L_k\big(y[k] - C\hat{\mathbf{x}}_k^-\big),\\
	P_k &= (I - L_kC)P_k^-.
\end{align}
\paragraph{Prediksi.}
\begin{align}
	\hat{\mathbf{x}}_{k+1}^- &= \Phi\hat{\mathbf{x}}_k + \Gamma u[k],\\
	P_{k+1}^- &= \Phi P_k\Phi^\top + \Gamma_wQ_n\Gamma_w^\top.
\end{align}
Matriks $P$ adalah kovarians galat estimasi. Untuk sistem LTI, $P_k^-$ dan $L_k$ konvergen ke nilai tunak (solusi DARE dual), sehingga mikrokontroler cukup menyimpan $L$ konstan. Intuisinya: bila $R_n$ besar (sensor buruk), filter lebih percaya model dan $L$ kecil; bila $Q_n$ besar (model buruk), filter lebih percaya sensor dan $L$ besar.

\begin{contoh}[Filter Kalman motor posisi]
	Dengan derau proses pada masukan $Q_n = 0{,}5$ dan derau potensiometer $\sigma_v = 0{,}01$ rad, penguatan tunaknya $L_\infty = [0{,}765\ \ 5{,}326]^\top$ dengan kutub galat berjari-jari 0,44. Simulasi dengan hanya posisi yang diukur menghasilkan estimasi kecepatan yang mengikuti nilai sebenarnya, termasuk fluktuasi akibat derau proses, dengan RMS galat sekitar 0,11 rad/s. Penalaan praktis dimulai dengan $R_n = \sigma_v^2$ dari data sensor saat diam, lalu $Q_n$ ditala hingga estimasi cukup responsif tanpa terlalu berderau.
\end{contoh}

\section{LQG dan MPC}
\term{LQG}{linear quadratic Gaussian} menggabungkan LQR dengan filter Kalman: $u = -K_{\text{LQR}}\hat{\mathbf{x}}$. Prinsip pemisahan berlaku sehingga $K$ dan $L$ dirancang terpisah, dan hasilnya optimal untuk biaya kuadratik dengan derau Gaussian. Namun LQG tidak memiliki jaminan margin ketegaran, sehingga $D(z)G(z)$ selalu perlu diperiksa.

\term{\textit{Model predictive control}}{MPC} menyelesaikan masalah optimisasi berhorison $N$ setiap cacah:
\begin{equation}
	\min_{u_0,\dots,u_{N-1}}\sum_{i=0}^{N-1}\big(\mathbf{x}_i^\top Q\mathbf{x}_i + u_i^\top Ru_i\big) + \mathbf{x}_N^\top S\mathbf{x}_N
\end{equation}
dengan kendala $\mathbf{x}_{i+1} = \Phi\mathbf{x}_i + \Gamma u_i$, $\mathbf{x}_0 = \hat{\mathbf{x}}[k]$, $u_{\min} \le u_i \le u_{\max}$, dan batasan keadaan. Hanya $u_0$ yang diterapkan, lalu masalah diulang pada cacah berikutnya (\textit{receding horizon}). Kelebihan utamanya adalah penanganan batasan secara eksplisit. Tanpa batasan aktif dan dengan $S = P_{\text{DARE}}$, MPC identik dengan LQR. MPC memerlukan solver \textit{quadratic programming} setiap cacah, yang menantang untuk Arduino UNO tetapi layak pada mikrokontroler yang lebih kuat atau dengan horison pendek.

\begin{ringkasan}
	LQR: $K = (R + \Gamma^\top P\Gamma)^{-1}\Gamma^\top P\Phi$ dengan $P$ dari DARE. Aturan Bryson memberi bobot awal dan $R$ menggeser kompromi kinerja--usaha. Filter Kalman adalah pengamat \textit{current} dengan penguatan optimal dari $Q_n$ dan $R_n$. LQG menggabungkan keduanya; MPC menambahkan batasan.\\[3pt]
	\en{LQR minimises a quadratic cost via the discrete Riccati equation; Bryson's rule gives initial weights; the Kalman filter is the optimal current observer; LQG combines both; MPC adds explicit constraints through receding-horizon optimisation.}
\end{ringkasan}

\paragraph{Istilah kunci.} indeks kinerja kuadratik \en{(quadratic cost)}, persamaan Riccati \en{(Riccati equation)}, aturan Bryson \en{(Bryson's rule)}, filter Kalman \en{(Kalman filter)}, kovarians \en{(covariance)}, LQG, \textit{receding horizon}.

\section*{Latihan}
\begin{enumerate}
	\item Implementasikan rekursi Riccati mundur dan tunjukkan konvergensinya ke hasil \texttt{dlqr}.
	\item Ulangi contoh LQR dengan $Q = \text{diag}(1, 0{,}1)$. Apa efek bobot kecepatan?
	\item Tala $Q_n$ filter Kalman untuk data posisi Modul 3 dan bandingkan estimasi kecepatan dengan tachometer.
	\item Rancang LQG untuk motor posisi dan hitung margin fasanya.
	\item Jelaskan mengapa MPC tanpa batasan aktif sama dengan LQR bila bobot terminalnya solusi DARE.
\end{enumerate}
